Persistent-memory poisoning---an attacker-influenced false or harmful
claim written into an LLM agent's long-term memory---is an established
cross-session security problem distinct from single-turn prompt
injection, because the compromised state outlives the triggering
interaction and can be \emph{derived} into further memories through
summarization, paraphrase, or continuation. Prior work asks either
whether a poisoned label survives to gate a future action, or how to
rank suspicious memories after harm is already observed. We study a
different, operationally realistic moment: once a memory or source is
\emph{known} to be compromised---independent
of whether it has caused visible harm yet---how accurately can execution
provenance reconstruct its downstream blast radius in a \textbf{branching}
derivation graph, and what does containment cost, as a function of
retrieval fan-out, derivation depth, and attribution policy?

We build a generation harness that produces branching derivation
graphs under controlled conditions across three LLMs (gpt-4o-mini,
gpt-4o, Llama-3.3-70B-Instruct), 30 scenarios (24 adversarial across
four poison-injection styles, 6 clean controls), yielding 21,570
generated traces and 216,319 labeled derived memories (Cohen's
$\kappa = 0.87$--$0.88$ against blind human annotation). We characterize
a precision/recall frontier, not a single failure mode: conservative
provenance preserves near-complete recall but produces growing
false-positive inflation with retrieval fan-out and derivation depth
(H1); attribution thresholding reduces that inflation at a measurable
recall cost, pruning exactly the weakly-attributed edges that are
sometimes the only path to real contamination (H2); a harmful claim's
literal surface form survives rewording more reliably than a naive
marker-matching check would suggest, but not perfectly---a small
(0.79\%), highly non-uniform, and model-dependent laundering rate
persists at full corpus scale (H3); and a depth-aware containment
policy reduces over-quarantine by roughly 30\% relative to a flat
conservative policy, at a quantifiable and cross-model-uneven
missed-contamination cost (H4). A held-out validation slice of 20 real
public documents (240 traces, 2{,}400 derived memories) reproduces the
main containment tradeoff---depth-aware propagation cutting inflation
from 3.57 to 2.50, structural precision/recall consistent with the
synthetic corpus---but shows a substantially higher and more
heterogeneous surface-laundering rate (19.1\% [6.9\%, 33.1\%],
scenario-level), indicating the provenance/containment frontier
generalizes more robustly than the absolute laundering rate does. No
single model is uniformly safer across all three failure modes we
measure, a result an operator needs, not an artifact to average away.
